%===============================================================================
\section{Experimental Setup}\label{sec:setup}
%===============================================================================

\subsection{Data sets and rule-based field roles}\label{sec:data}
Table~\ref{tab:data} summarizes the four data sets. In all of them, candidate fields are
information that the disclosure rules~\cite{nea2024disclosure} release \emph{before} the
target, and targets are either specific information (unit output, Art.~17(3)), quantities not on
the disclosure list at hourly resolution (hourly line flow; only daily averages of important
lines are listed), or actual operating quantities released on the following day (actual load,
total generation, tie-line exchange; Appendix items 6.57--6.66). Using the release timing to
assign roles does not involve any temporal split of the data; all splits are random.

\emph{RTS-GMLC.} The IEEE Reliability Test System 2019 update~\cite{barrows2020rts} provides a
73-bus network with one year of realistic hourly load, wind, solar and hydro series. We clear
an hourly DC optimal power flow with piecewise-linear offers built from the published heat-rate
curves (8,784 hours, 26.7\% congested). Candidates are the 23 fields a market would publish
early (zonal load forecasts, renewable and hydro forecasts, reserve requirements, system energy
price and nine nodal prices, three congestion indicators); targets are the outputs of a coal and
a combined-cycle unit and the hourly flow on internal line C35.

\emph{NEM.} One year (2024) of the Australian National Electricity Market~\cite{aemo2024nemweb},
averaged to hours. Candidates are regional demand, available capacity, semi-scheduled wind and
solar forecasts, regional prices and interconnector flows (30 fields); targets are the actual
outputs of a coal unit (BW01), a hydro unit (TUMUT3) and a gas combined-cycle unit (PPCCGT).
Unit output is public in the NEM but is specific information under the Chinese rules; this data set supplies observed unit outputs for the stated disclosure scenario.

\emph{PJM and CAISO.} System-level hourly data for one year. Targets are the actual quantities
released on the next day: PJM metered load, total generation and net actual interchange, and
CAISO actual load. Candidates are forecasts, day-ahead and real-time prices and their
components, ancillary-service requirements and cleared volumes, and schedules (22--35 fields).
Fields released on the same day as the target (fuel-mix generation, actual renewable output,
other balancing areas' actual load) are excluded, as are the load forecasts of the target area,
which are ex-ante copies of the target (Section~\ref{sec:problem}). Near-copies that are distinct
quantities (e.g.\ sub-area forecasts for CAISO load) are kept; removing them is analysed as a
sensitivity case.

\emph{Controlled mechanism.} A fifth data set is built on the RTS-GMLC market with a private hourly
offer markup of one unit as the target and eleven fields chosen to span the known mechanism
(Section~\ref{sec:res-mech}); all 2,047 of its field sets are retrained.

\begin{table}[!tb]
\centering\small
\caption{Audit data and reference coverage. Rows give training/test sample sizes. Each reference table contains all nonempty field sets of size at most four. The field-role basis is described in the text and Table~\ref{tab:fields}.}\label{tab:data}
\setlength{\tabcolsep}{5pt}
\begin{tabular}{llrrr}
\toprule
Data & Targets & Train/test & Fields & Sets\\
\midrule
RTS-GMLC & Units 223, 221; line C35 & 6149/2635 & 23 & 10,902\\
NEM 2024 & BW01, TUMUT3, PPCCGT & 6149/2635 & 30 & 31,930\\
PJM load & Metered load & 4589/1967 & 22 & 9,108\\
PJM gen./interch. & Generation; net interchange & 4589/1967 & 24 & 12,950\\
CAISO & Actual load & 6105/2617 & 35 & 59,535\\
\bottomrule
\end{tabular}
\end{table}

\subsection{Ground truth}
All variables are standardized with training statistics; 70/30 random train/test splits are
used, with 15\% of the training rows as validation. The attacker family has three members: a
single-target multilayer perceptron (two hidden layers of 128 units, dropout 0.15, Adam,
early stopping on validation), a multi-target perceptron trained on all targets of a data set,
and histogram gradient boosting with two configurations. For every set of up to four fields,
every member is trained on that set alone; the member and epoch are chosen on validation and the
test $\Rtwo$ is reported, followed by the monotone closure~\eqref{eq:closure}. This amounts to
124,425 field sets and several hundred thousand trained models. These exhaustive tables provide the reference for evaluating the estimators and field scores.

\subsection{Baselines}
Field scores compared with $M^{(2)}$ are: Pearson $r^2$, squared Spearman correlation, mutual
information~\cite{kraskov2004mi}, distance correlation~\cite{szekely2007dcor}, single-field
retraining $M^{(0)}$ (current practice), LOCO~\cite{lei2018loco}, permutation
importance~\cite{breiman2001rf}, SAGE with marginal imputation~\cite{covert2020sage}, and an
inference-graph score that follows the propagation idea of~\cite{hu2026ignn} without its expert
inputs: edge weights are the single-field predictability of one field from another (gradient
boosting, validation $\Rtwo$), the target node carries risk one, and risk is propagated over three
hops with decay 0.5. All baseline scores are computed on training rows only.

Estimators compared with ours avoid retraining the attacker family for each set in different ways:
per-set linear and quadratic regression (retraining the simplest attackers); a model trained on all fields with the removed
fields replaced by their mean~\cite{williamson2020spvim}, linearized around its parameters with a ridge correction per set
(LazyVI~\cite{gao2022lazyvi}), or warm-started and trained with early stopping per set~\cite{sun2025warmstart}; a
surrogate trained under random masks whose output is read directly~\cite{covert2021explaining,jethani2022fastshap}; and the
tabular foundation model TabPFN v2~\cite{hollmann2025tabpfn} used in context. Design choices of our estimator---the feature
map, the use of seeds, the safeguards and the pre-training masks---are compared separately (Table~\ref{tab:design}); details
of all implementations are given in Appendix~\ref{app:impl}.

\subsection{Metrics}
For set values: mean absolute error of $\hat V$. For field risk: mean absolute error and bias of
$\hat M^{(2)}$, Spearman correlation across fields, and the certified ratio
$\sum_i L^{(2)}_{i\to y}/\sum_i M^{(2)}_{i\to y}$ when $k$ backgrounds per field are retrained.
For decisions at thresholds $\tau\in\{0.5,0.7\}$: recall and precision of critical fields; the
number of fields withheld and the number of minimal unsafe sets of size $\le 3$ left unbroken.
